\section{Conclusiones}

\subsection{Hallazgos principales}

El trabajo modeló tres series de naturaleza distinta y obtuvo resultados de calidad desigual,
circunstancia que conviene exponer sin uniformar.

Las dos series de tráfico presentan estacionalidad diaria inequívoca, con período de
veinticuatro horas, verificada tanto por la inspección del perfil horario como por los picos de
la función de autocorrelación en los rezagos 24, 48 y 72. Los modelos SARIMA seleccionados
alcanzan errores porcentuales del 0,79 y del 3,90 por ciento sobre un horizonte de 48 horas,
magnitudes aceptables para una serie de tráfico operativo.

Ese desempeño, sin embargo, convive con un diagnóstico de residuos desfavorable. La prueba de
Ljung-Box rechaza la incorrelación a partir del rezago 24 en la serie del balanceador, y en
todos los rezagos examinados en la serie de punto de venta. Ambos modelos dejan estructura sin
explicar, concentrada en los rezagos estacionales. Los pronósticos merecen, por lo tanto, menor
confianza que la sugerida por sus métricas de error.

La serie de altitud presenta el comportamiento opuesto. Sin componente estacional, es un proceso integrado que una sola diferencia regular vuelve
estacionario. El modelo
ARIMA$(1,1,3)$ exhibe parámetros significativos tanto individual como conjuntamente, y sus
residuos no se distinguen del ruido blanco en ningún rezago examinado. Es el único de los tres que supera el diagnóstico sin reservas.

El modelo multivariado, construido sobre las dos series de tráfico con dos rezagos, arroja una
relación de precedencia unidireccional: los rezagos de la serie de punto de venta contribuyen a
explicar la del balanceador, sin reciprocidad.

\subsection{Dos lecciones metodológicas}

Dos episodios del trabajo trascienden el resultado particular.

El primero es la selección del orden en el modelo multivariado. El criterio de Akaike sugería
24 rezagos, esto es, 98 parámetros estimados a partir de 241 observaciones. La especificación
resultante predecía peor que una de diez parámetros. Los criterios de información orientan, no dictaminan, y ante discrepancia entre ellos el desempeño fuera de muestra es el árbitro más confiable. La misma tensión reapareció, con menor intensidad, al
seleccionar el modelo estacional de la serie de punto de venta.

El segundo es de orden más elemental: una métrica mal elegida puede invertir por completo la
lectura de un resultado. El error porcentual absoluto medio del 97,20 por ciento que arrojaba
el modelo multivariado no medía su calidad sino la inadecuación de aplicar una métrica relativa
sobre una serie que oscila alrededor de cero. Expresado en la escala original, el mismo modelo
alcanza un error del 1,31 por ciento.

\subsection{Límites del trabajo}

\begin{enumerate}
\item \textbf{Tamaño muestral.} Las series de tráfico abarcan aproximadamente un mes, y la de
altitud, 457 observaciones. Ninguna admite validación mediante ventanas móviles ni evaluación
sobre múltiples horizontes independientes.

\item \textbf{Estacionalidad semanal no modelada.} El mapa de calor revela que la serie del
balanceador presenta niveles sensiblemente inferiores los días miércoles y jueves. Un período
muestral de un mes ofrece apenas cuatro repeticiones de cada día de la semana, cantidad
insuficiente para estimar una componente estacional de período 168. La autocorrelación residual detectada en el diagnóstico encaja con esa omisión.

\item \textbf{Solapamiento limitado entre las series de tráfico.} La intersección temporal
abarca trece días y 289 observaciones tras la transformación. Toda conclusión del modelo
multivariado queda condicionada por esa brevedad.

\item \textbf{Calendarios disjuntos.} La serie de altitud no participa del modelo multivariado
por corresponder a un período que no solapa con el de las restantes. La limitación es
estructural y no admite solución dentro de los datos disponibles.

\item \textbf{Alcance de la causalidad de Granger.} El contraste establece precedencia temporal
con contenido predictivo. No autoriza a afirmar que una serie cause a la otra en sentido
estructural.

\item \textbf{Un único vuelo.} Los resultados de la serie de altitud describen ese vuelo. No se
generalizan a otros vuelos ni a otras aeronaves.

\item \textbf{Normalidad de los residuos.} El contraste de Jarque-Bera rechaza la normalidad en
el modelo de la serie de altitud. Los intervalos de predicción nominales son, por eso, aproximados.
\end{enumerate}

\subsection{Líneas de trabajo futuro}

La limitación de mayor peso es la brevedad del período muestral, y corregirla destrabaría las
restantes. Un año de registro permitiría estimar la componente semanal cuya ausencia se
manifiesta en la autocorrelación residual, y ampliaría el solapamiento entre las series de
tráfico hasta un tamaño donde el modelo multivariado admitiera órdenes superiores sin
comprometer la fiabilidad.

En segundo término, el modelo del balanceador se puede simplificar. Los
parámetros de la parte regular carecen de significatividad, de modo que una especificación
$(0,1,0)\times(1,1,1)_{24}$ alcanzaría un ajuste equivalente con dos parámetros menos.

Por último, la serie de altitud ofrece una extensión natural que este trabajo no exploró: la telemetría de vuelo registra en paralelo la velocidad aerodinámica junto con la
altitud, y ambas magnitudes están vinculadas por el intercambio entre energía potencial y
cinética. Un modelo multivariado sobre ese par contaría con series perfectamente alineadas y
con una interpretación física directa de la función impulso-respuesta.
